\newpage
\section{Test 12: Fine-Tuned Qwen3-4B BERT Evaluation}

\noindent\textbf{Experiment Objective}

We evaluated the merged fine-tuned Qwen3-4B model on 10,000 PubMedQA examples and compared its answers with the dataset's reference answers using BERTScore. We also compared the results with the base model evaluated using the updated prompt in Test 11.

\subsection{Experimental Configuration}

\begin{itemize}
    \item \textbf{Base Model:} \texttt{Qwen/Qwen3-4B-Base}
    \item \textbf{Fine-Tuned Model:} \texttt{merged\_adapters\_new2}
    \item \textbf{Fine-Tuning Method:} LoRA, rank 16 and alpha 16
    \item \textbf{Training Data:} First 10,000 examples from the PubMedQA artificial training split, divided into 80\% training and 20\% validation examples
    \item \textbf{Training Configuration:} One epoch, learning rate $10^{-3}$, and maximum sequence length of 2,048 tokens
    \item \textbf{Inference Script:} \texttt{vllm\_fine.py}
    \item \textbf{Generation:} Temperature 0 and a maximum of 400 generated tokens
    \item \textbf{Evaluation Samples:} 10,000
    \item \textbf{BERTScore Model:} \texttt{microsoft/deberta-xlarge-mnli}
\end{itemize}

\subsection{BERTScore Results}

The base and fine-tuned result files contain the same 10,000 reference answers in the same row order. The following table compares their scores. The 50th percentile is the median.

\begin{center}
\begin{tabular}{lrrr}
\hline
\textbf{Metric} & \textbf{Base} & \textbf{Fine-Tuned} & \textbf{Change} \\
\hline
Mean precision & 0.6308 & 0.3472 & -0.2836 \\
Mean recall & 0.6604 & 0.4868 & -0.1736 \\
Mean F1 & 0.6436 & 0.3996 & -0.2440 \\
Median F1 (50th percentile) & 0.6419 & 0.3974 & -0.2445 \\
Standard deviation of F1 & 0.0565 & 0.0623 & +0.0058 \\
Minimum F1 & 0.2798 & 0.1720 & -0.1078 \\
10th percentile F1 & 0.5807 & 0.3249 & -0.2558 \\
25th percentile F1 & 0.6095 & 0.3594 & -0.2501 \\
75th percentile F1 & 0.6768 & 0.4399 & -0.2369 \\
90th percentile F1 & 0.7128 & 0.4812 & -0.2316 \\
Maximum F1 & 0.9300 & 0.6110 & -0.3190 \\
\hline
\end{tabular}
\end{center}

The fine-tuned model had a lower F1 score on 99.54\% of the paired examples. Its mean output length was approximately 1,749 characters, compared with 326 characters for the base model. In addition, 26.9\% of the fine-tuned outputs repeated the evaluation instruction beginning with ``Write one or two concise declarative sentences,'' and 52.4\% exceeded 1,800 characters. These lengths are measured in characters, not tokens.

\subsection{Examples of Fine-Tuned Model Outputs}

The excerpts below show text generated by the fine-tuned model. Ellipses indicate omitted text.

\begin{quote}
\textbf{Example 1: Repeated answers}\\
Question: Is the variation in high sensitive cardiac troponin concentration during haemodialysis treatment similar to the biological variation observed in stable end stage renal disease patients?\\
Output: 1

The variation in high sensitive cardiac troponin concentration during haemodialysis treatment is similar to the biological variation observed in stable end stage renal disease patients.

2

3

The variation in high sensitive cardiac troponin concentration during haemodialysis treatment is similar to the biological variation observed in stable end stage renal disease patients. ...\\
Reference: The CV\\
F1: 0.1720
\end{quote}

\begin{quote}
\textbf{Example 2: Repeated numbering without an answer}\\
Output: 1

Answer: 2

Answer: 3

Answer: 4

Answer: 5 ...\\
Reference: As sub group sizes were small, data suggesting a trait-like deficit in ToM should be interpreted with caution.\\
F1: 0.2183
\end{quote}

\begin{quote}
\textbf{Example 3: Repeated and expanding claims}\\
Question: Does inflamed skin predispose to sensitization to less potent allergens?\\
Output: 1. Inflamed skin predisposes to sensitization to less potent allergens. 2. Inflamed skin is a risk factor for sensitization to less potent allergens, and this may be due to the increased expression of the TLR2 receptor in inflamed skin. 3. Inflamed skin is a risk factor for sensitization to less potent allergens, and this may be due to the increased expression of the TLR2 receptor in inflamed skin, which may be due to the increased expression of the TLR2 receptor in inflamed skin. ...\\
Reference: Association cannot prove causation.\\
F1: 0.1852
\end{quote}

\subsection{Interpretation and Possible Causes}

The fine-tuned model often continues generating numbered items or repeats text instead of ending after a concise answer. The repeated and unrelated text is consistent with the large drop in BERTScore precision. BERTScore measures similarity to the reference; it does not by itself establish whether an answer is medically correct. Some references are short fragments or limitations, so individual scores should be interpreted alongside the question, study context, and generated text.

The training and inference inputs differ. Training used system and user messages containing the question and study context, with the reference answer as the completion. Evaluation instead used a plain-text instruction ending in ``Answer:'' and set the context variable to false. The \texttt{--context} option is defined in the inference script, but the script currently sets \texttt{context = False} regardless. This difference in both message format and available information is a strong possible cause of the poor output; it has not yet been isolated experimentally.

A second possible cause is stopping behavior. The saved chat template ends assistant messages with \texttt{<|im\_end|>} (token 151645), while the merged model's configured EOS token is \texttt{<|endoftext|>} (token 151643). The existing outputs do not include generated token IDs or finish reasons, so we cannot establish whether the model generated an unrecognized end-of-message token or reached the 400-token limit. This token mismatch should be checked directly.

The learning rate of $10^{-3}$ may also have changed the model too aggressively. The saved trainer state reports evaluation loss 1.8238 for checkpoint 223; that loss alone does not measure the quality of generated answers. The merge script defaults to checkpoint 223 and an output directory named \texttt{merged\_adapters\_new}, whereas inference loads \texttt{merged\_adapters\_new2}. The merge command should be checked to confirm the adapter used for the evaluated model.

\subsection{Next Debugging Steps}

\begin{enumerate}
    \item Evaluate the base and fine-tuned models on the same held-out examples with the same message format, study context, and generation settings. The current evaluation samples from the dataset's training split, so it may include examples used during fine-tuning.
    \item Save generated token IDs, output token counts, and finish reasons. Check for token 151645 and compare a run that explicitly stops at \texttt{<|im\_end|>} with the existing stopping settings.
    \item Compare output length, repeated text, instruction echoes, BERTScore precision, recall, and F1, and inspect a random sample as well as low-scoring examples. Try a shorter generation limit as a separate diagnostic.
    \item After aligning prompts and the evaluation split, compare checkpoints and train a controlled run with a lower learning rate, such as $10^{-4}$, while keeping other settings fixed.
\end{enumerate}
